% Options for packages loaded elsewhere
\PassOptionsToPackage{unicode}{hyperref}
\PassOptionsToPackage{hyphens}{url}
\documentclass[
  ignorenonframetext,
  aspectratio=169,
]{beamer}
\newif\ifbibliography
\usepackage{pgfpages}
\setbeamertemplate{caption}[numbered]
\setbeamertemplate{caption label separator}{: }
\setbeamercolor{caption name}{fg=normal text.fg}
\beamertemplatenavigationsymbolsempty
% remove section numbering
\setbeamertemplate{part page}{
  \centering
  \begin{beamercolorbox}[sep=16pt,center]{part title}
    \usebeamerfont{part title}\insertpart\par
  \end{beamercolorbox}
}
\setbeamertemplate{section page}{
  \centering
  \begin{beamercolorbox}[sep=12pt,center]{section title}
    \usebeamerfont{section title}\insertsection\par
  \end{beamercolorbox}
}
\setbeamertemplate{subsection page}{
  \centering
  \begin{beamercolorbox}[sep=8pt,center]{subsection title}
    \usebeamerfont{subsection title}\insertsubsection\par
  \end{beamercolorbox}
}
\usetheme[]{Madrid}
\usefonttheme{serif} % use mainfont rather than sansfont for slide text
% Prevent slide breaks in the middle of a paragraph
\widowpenalties 1 10000
\raggedbottom
\AtBeginPart{
  \frame{\partpage}
}
\AtBeginSection{
  \ifbibliography
  \else
    \frame{\sectionpage}
  \fi
}
\AtBeginSubsection{
  \frame{\subsectionpage}
}
\usepackage{iftex}
\ifPDFTeX
  \usepackage[T1]{fontenc}
  \usepackage[utf8]{inputenc}
  \usepackage{textcomp} % provide euro and other symbols
\else % if luatex or xetex
  \usepackage{unicode-math} % this also loads fontspec
  \defaultfontfeatures{Scale=MatchLowercase}
  \defaultfontfeatures[\rmfamily]{Ligatures=TeX,Scale=1}
\fi
\ifPDFTeX\else
  % xetex/luatex font selection
    \setmainfont[]{Latin Modern Roman}
    \setsansfont[]{Latin Modern Sans}
    \setmonofont[]{Latin Modern Mono}
\fi
% Use upquote if available, for straight quotes in verbatim environments
\IfFileExists{upquote.sty}{\usepackage{upquote}}{}
\IfFileExists{microtype.sty}{% use microtype if available
  \usepackage[]{microtype}
  \UseMicrotypeSet[protrusion]{basicmath} % disable protrusion for tt fonts
}{}
\makeatletter
\@ifundefined{KOMAClassName}{% if non-KOMA class
  \IfFileExists{parskip.sty}{%
    \usepackage{parskip}
  }{% else
    \setlength{\parindent}{0pt}
    \setlength{\parskip}{6pt plus 2pt minus 1pt}}
}{% if KOMA class
  \KOMAoptions{parskip=half}}
\makeatother
\usepackage{longtable,booktabs,array}
\usepackage{calc} % for calculating minipage widths
\usepackage{caption}
% Make caption package work with longtable
\makeatletter
\def\fnum@table{\tablename~\thetable}
\makeatother
\usepackage{graphicx}
\makeatletter
\newsavebox\pandoc@box
\newcommand*\pandocbounded[1]{% scales image to fit in text height/width
  \sbox\pandoc@box{#1}%
  \Gscale@div\@tempa{\textheight}{\dimexpr\ht\pandoc@box+\dp\pandoc@box\relax}%
  \Gscale@div\@tempb{\linewidth}{\wd\pandoc@box}%
  \ifdim\@tempb\p@<\@tempa\p@\let\@tempa\@tempb\fi% select the smaller of both
  \ifdim\@tempa\p@<\p@\scalebox{\@tempa}{\usebox\pandoc@box}%
  \else\usebox{\pandoc@box}%
  \fi%
}
% Set default figure placement to htbp
\def\fps@figure{htbp}
\makeatother
\setlength{\emergencystretch}{3em} % prevent overfull lines
\providecommand{\tightlist}{%
  \setlength{\itemsep}{0pt}\setlength{\parskip}{0pt}}
\usepackage{bookmark}
\IfFileExists{xurl.sty}{\usepackage{xurl}}{} % add URL line breaks if available
\urlstyle{same}
\hypersetup{
  pdftitle={Rice Growth-Stage, Cropping Intensity \& Production from S1+S2 Fusion},
  pdfauthor={Dr.~Firman Hadi},
  hidelinks,
  pdfcreator={LaTeX via pandoc}}

\title{Rice Growth-Stage, Cropping Intensity \& Production from S1+S2
Fusion}
\subtitle{MOGPR Sentinel-1 + Sentinel-2 --- Conclusions (Java,
2024-2026)}
\author{Dr.~Firman Hadi}
\date{June 2026}

\begin{document}
\frame{\titlepage}

\begin{frame}{The question}
\phantomsection\label{the-question}
\textbf{Can we map rice growth stage --- and from it cropping intensity
and production --- better than with Sentinel-1 alone?}

\begin{itemize}
\tightlist
\item
  The TA-2026 reports built a strong \textbf{S1-only} paddy map (95\%
  CV, \textasciitilde3.0 M ha).
\item
  But their \textbf{growth-stage} model was \emph{not validated} (32\%
  vs field data) and they parked \textbf{S1 + S2 fusion} as long-term
  future work.
\item
  This work builds and validates that fusion.
\end{itemize}
\end{frame}

\begin{frame}{Approach}
\phantomsection\label{approach}
\begin{itemize}
\tightlist
\item
  \textbf{MOGPR} fuses Sentinel-1 VH (SAR, all-weather) + Sentinel-2
  NDVI (optical) into one gap-filled curve per pixel.
\item
  MOGPR is a \textbf{feature generator}, not the final model.
\item
  A trained \textbf{MiniROCKET + LightGBM} classifier reads the fused
  curve -\textgreater{} 6 / 3 growth phases.
\item
  Production = area in generative phase (4+5) x 5.8 t/ha, per 12-day
  period, Java-wide, 50 m.
\end{itemize}
\end{frame}

\begin{frame}{Finding 1 --- Fusion beats SAR-only for phase}
\phantomsection\label{finding-1-fusion-beats-sar-only-for-phase}
Generative-phase F1 (leave-one-region-out):

\begin{longtable}[]{@{}ll@{}}
\toprule\noalign{}
Method & F1 \\
\midrule\noalign{}
\endhead
Rule-based on SAR curve & \textasciitilde22\% \\
\textbf{VH-only} & \textbf{52\%} \\
\textbf{VH + MOGPR optical} & \textbf{67-68\%} \\
idmai VH-CNN (in-sample) & 64\% \\
\bottomrule\noalign{}
\end{longtable}

\begin{itemize}
\tightlist
\item
  \textbf{Optical is the dominant signal} for phase; SAR alone cannot
  carry it.
\item
  This is the validated phase model the reports lacked.
\end{itemize}
\end{frame}

\begin{frame}{Finding 2 --- Multi-season training is non-negotiable}
\phantomsection\label{finding-2-multi-season-training-is-non-negotiable}
A model trained only on the wet season \textbf{collapses} on the dry
season:

\begin{longtable}[]{@{}lll@{}}
\toprule\noalign{}
Training & Dry-season accuracy & Dry-season generative F1 \\
\midrule\noalign{}
\endhead
Wet-only & 36.9\% & 60.7\% \\
\textbf{Wet + dry (multi-season)} & \textbf{78.6\%} & \textbf{91.6\%} \\
\bottomrule\noalign{}
\end{longtable}

\begin{itemize}
\tightlist
\item
  Fixed using existing dry-season field labels (no new data needed).
\item
  Independently reproduces the S1-only literature's central finding ---
  on the fused signal.
\end{itemize}
\end{frame}

\begin{frame}{Finding 3 --- Production (method demonstration)}
\phantomsection\label{finding-3-production-method-demonstration}
\begin{longtable}[]{@{}lll@{}}
\toprule\noalign{}
Year & Cropping intensity & Production \\
\midrule\noalign{}
\endhead
2024 & 2.18 & \textbf{28.8 Mt} \\
2025 & 2.22 & \textbf{29.4 Mt} \\
2024-2025 combined & 2.22 & \textbf{\textasciitilde29.5 Mt/yr} \\
\bottomrule\noalign{}
\end{longtable}

\begin{itemize}
\tightlist
\item
  Harvest area \textasciitilde5.0 M ha \textbf{= close to BPS luas panen
  (\textasciitilde5.34 M)} --- an independent sanity check.
\item
  Detected-paddy denominator (2.28 M ha), flat 5.8 t/ha.
\end{itemize}
\end{frame}

\begin{frame}{Finding 4 --- A real interannual signal}
\phantomsection\label{finding-4-a-real-interannual-signal}
\begin{figure}
\centering
\includegraphics[width=0.85\linewidth,height=\textheight,keepaspectratio]{output/production/ci3_highlight_2024_2025.png}
\caption{Triple-cropping (IP=3) concentrates in the Bengawan Solo
corridor}
\end{figure}

\begin{itemize}
\tightlist
\item
  2025 \textgreater{} 2024 (wetter, post-El-Nino) --- confirmed real,
  not artefact.
\item
  \textasciitilde Half of Java's triple-cropping sits in the
  \textbf{Bengawan Solo} basin (water-secure year-round).
\end{itemize}
\end{frame}

\begin{frame}{Java cropping intensity --- combined 2024-2025}
\phantomsection\label{java-cropping-intensity-combined-2024-2025}
\begin{figure}
\centering
\includegraphics[width=0.95\linewidth,height=\textheight,keepaspectratio]{output/production/java_combined_2024_2025/java_cropping_intensity_combined.png}
\caption{Annual cropping intensity, 24-month combined count}
\end{figure}
\end{frame}

\begin{frame}{Honest limits}
\phantomsection\label{honest-limits}
\begin{itemize}
\tightlist
\item
  \textbf{Method demonstration, not official numbers:} detected-mask
  denominator + flat yield.
\item
  Accuracy is \textbf{optimistic} --- validated on a random (not
  spatial) split, same-campaign labels.
\item
  6-class is weak (\textasciitilde51\%); robust products are
  \textbf{3-class} and the \textbf{generative} class.
\end{itemize}
\end{frame}

\begin{frame}{Conclusion}
\phantomsection\label{conclusion}
\begin{itemize}
\tightlist
\item
  \textbf{S1 + S2 MOGPR fusion materially improves rice growth-stage
  mapping over S1-only}, and turns it into a working
  \textbf{cropping-intensity and production} product for Java.
\item
  It delivers --- ahead of schedule --- the fusion the program had
  parked as future work.
\item
  The S1 binary-paddy results stand; this \textbf{adds} a validated
  phase -\textgreater{} IP -\textgreater{} production capability.
\end{itemize}
\end{frame}

\begin{frame}{Next step --- the decisive one}
\phantomsection\label{next-step-the-decisive-one}
\textbf{TA-2027 multi-DI field validation} (Klambu, Rentang, Jatiluhur,
Semarang, E. Java, rainfed):

\begin{itemize}
\tightlist
\item
  Spatial leave-one-DI-out test (not random) + a dry-season campaign.
\item
  Validate IP against tracked actual harvests; production vs KSA/BPS.
\item
  Converts ``promising demonstration'' -\textgreater{} ``operationally
  trustworthy.''
\item
  \textbf{Pilot already prepared:} 100 lag-adjusted Semarang points
  awaiting survey.
\end{itemize}
\end{frame}

\end{document}
